\section{Algoritmi}

Riportiamo in pseudocodice le due operazioni fondamentali su un blocco RowMask23: la ricostruzione della matrice e il prodotto diretto con un vettore. Le coordinate seguono la disposizione della \autoref{sec:layout-coordinate}: diagonale in $c[0..127]$, base $b_s$ in $c[127+s]$, correzioni $e_{i,k}$ in $c[255+23(i-1)+k]$.

\begin{algorithm}[H]
\caption{Ricostruzione di un blocco RowMask23}
\KwIn{coordinate $c\in\mathbb{R}^{3176}$, scala $\alpha$}
\KwOut{matrice $W\in\mathbb{R}^{128\times128}$}
\For{$i\leftarrow0$ \KwTo $127$}{
  \For{$j\leftarrow0$ \KwTo $127$}{
    $s\leftarrow i\oplus j$\;
    \lIf{$s=0$}{$W_{ij}\leftarrow c[i]$}
    \lElse{$W_{ij}\leftarrow c[127+s]$}
  }
}
\For{$i\leftarrow1$ \KwTo $127$}{
  \For{$k\leftarrow0$ \KwTo $22$}{
    $s\leftarrow1+((13k+7i)\bmod127)$\;
    $W_{i,\,i\oplus s}\leftarrow W_{i,\,i\oplus s}+c[255+23(i-1)+k]$\;
  }
}
\Return $\alpha W$\;
\end{algorithm}

\begin{algorithm}[H]
\caption{Prodotto diretto $y=Wx$ senza ricostruzione}
\KwIn{coordinate $c$, scala $\alpha$, vettore $x\in\mathbb{R}^{128}$}
\KwOut{$y=\alpha\,\Phi(c)\,x$}
$u_0\leftarrow0$; \lFor{$s\leftarrow1$ \KwTo $127$}{$u_s\leftarrow c[127+s]$}
$v\leftarrow x$\;
\For(\tcp*[f]{trasformata di Walsh--Hadamard}){$h\in\{1,2,4,\dots,64\}$}{
  \ForEach{$i$ con $i\land h=0$}{
    $(u_i,u_{i+h})\leftarrow(u_i+u_{i+h},\,u_i-u_{i+h})$\;
    $(v_i,v_{i+h})\leftarrow(v_i+v_{i+h},\,v_i-v_{i+h})$\;
  }
}
$u\leftarrow u\odot v$\;
\For(\tcp*[f]{trasformata inversa}){$h\in\{1,2,4,\dots,64\}$}{
  \ForEach{$i$ con $i\land h=0$}{
    $(u_i,u_{i+h})\leftarrow(u_i+u_{i+h},\,u_i-u_{i+h})$\;
  }
}
\For{$i\leftarrow0$ \KwTo $127$}{
  $y_i\leftarrow u_i/128 + c[i]\,x_i$\;
  \If{$i\ge1$}{
    \For{$k\leftarrow0$ \KwTo $22$}{
      $s\leftarrow1+((13k+7i)\bmod127)$\;
      $y_i\leftarrow y_i + c[255+23(i-1)+k]\;x_{i\oplus s}$\;
    }
  }
}
\Return $\alpha y$\;
\end{algorithm}

Nel kernel GPU i due cicli esterni di ogni stadio della trasformata sono eseguiti in parallelo dai $128$ thread del gruppo, con una barriera di sincronizzazione tra uno stadio e il successivo.

\section{Rubrica del giudizio narrativo}
\label{sec:appendice-rubrica}

Il giudice riceve, per ciascun identificativo, il prompt e la continuazione, senza alcuna informazione sul modello che l'ha prodotta. Deve leggere tutto il testo, valutarlo una sola volta e motivare il giudizio con un riferimento concreto al testo. Le istruzioni, tradotte dall'originale, sono le seguenti.

\paragraph{Grammatica (1--10).} 9--10: testo scorrevole e grammaticalmente corretto; 7--8: errori minori che non ostacolano la comprensione; 6: testo comprensibile nonostante errori localizzati; 4--5: errori ricorrenti che ostacolano la comprensione; 1--3: testo in gran parte incomprensibile.

\paragraph{Coerenza (1--10).} 9--10: identità, attributi, oggetti e successione degli eventi stabili; 7--8: ambiguità marginali o piccoli dettagli imprecisi; 6: nucleo narrativo interpretabile senza contraddizioni sostanziali; 4--5: contraddizioni non risolte su identità, proprietà o eventi che cambiano il senso di una scena; 1--3: identità o eventi centrali contraddittori o non interpretabili.

\paragraph{Trama (1--10).} 9--10: sviluppo chiaro e ben collegato; 7--8: semplice ma coerente; 6: trama minima comprensibile con eventi collegati; 4--5: sequenze scollegate, ripetitive o con causalità essenziale mancante; 1--3: assenza di una trama comprensibile.

\paragraph{Categoria.} S: tutti e tre i voti almeno 6 e nucleo narrativo coerente. P: nucleo riconoscibile ma almeno una criticità sostanziale, con almeno un voto non superiore a 5. N: perdita grave di coerenza o comprensibilità, con almeno un voto non superiore a 3.

\paragraph{Criteri aggiuntivi.} Non penalizzare in sé un finale negativo, la personificazione, gli animali parlanti, la magia, una trama semplice o un finale troncato dal limite di generazione. Una trasformazione esplicitamente motivata è ammessa; un cambio immotivato di identità o proprietà è un errore. Non è richiesta una morale. Un nuovo personaggio è ammesso se non sostituisce in modo ambiguo il protagonista. Va penalizzata ogni contraddizione oggettiva rispetto al prompt, indipendentemente dalla grammatica.

\paragraph{Esempi di ancoraggio.} ``Mia found a red ball. She gave it to Tom. Tom played with it.'': S. ``Mia gave the only ball to Tom, then said she had never given it away.'': P. ``Mia was a cat. Without explanation Mia was a boy, and Tom was Mia.'': N, se l'identità centrale diventa incomprensibile.

\section{Riproducibilità}

Tutti i numeri dei Capitoli~\ref{cap:capitoli/06-risultati} e~\ref{cap:capitoli/07-sistemi} sono prodotti da script conservati con i rispettivi risultati e con i log completi. La \autoref{tab:artefatti} elenca le impronte SHA-256 dei principali artefatti.

\begin{table}[H]
\centering
\small
\begin{tabular}{ll}
\toprule
Artefatto & SHA-256 (inizio) \\
\midrule
Pesi STERA 30M, checkpoint finale & \texttt{418becfbcaddcefd\dots} \\
Tokenizer & \texttt{f6ed3d307010c244\dots} \\
Prompt e racconti TEST300 & \texttt{95bcffdbc25502dd\dots} \\
TinyStories-33M & \texttt{41316d3bad3cf776\dots} \\
TinyStories-8M & \texttt{22c355bfabebc1f6\dots} \\
\bottomrule
\end{tabular}
\caption{Impronte dei principali artefatti sperimentali.}
\label{tab:artefatti}
\end{table}

Il repository della tesi contiene, per ogni misura, lo script che la produce e i dati grezzi. Le analisi delle componenti delle proiezioni usano lo script di valutazione con le opzioni per la quantizzazione del low-rank, la riduzione del rango e l'eliminazione o la quantizzazione di scale e guadagni; i giudizi ciechi delle varianti sono conservati con i rispettivi insiemi anonimizzati e abbinamenti. Le misure di tempo sono prodotte da script che eseguono le cinque campagne a sistema scarico e scrivono un riepilogo per ogni studio: confronto fra i due motori, confronto alla pari con i riferimenti nello stesso motore, lunghezza dell'ingresso, strategie di prefill e low-rank a 8 bit. I motori aggiunti per questi studi sono il motore denso per i riferimenti, la variante del motore compatto con low-rank a 8 bit e il prefill con espansione al bisogno, che contiene il kernel di espansione verificato bit per bit rispetto alle matrici dense.
